%% General definitions
%\documentclass{article} %% Determines the general format.
\documentclass[10pt,a4paper]{article}
\usepackage{a4wide} %% paper size: A4.
\usepackage[utf8]{inputenc} %% This file is written in UTF-8.
%% Some editors on Windows cannot save files in UTF-8.
%% If there is a problem with special characters not showing up
%% correctly, try switching "utf8" to "latin1" (ISO 8859-1).
\usepackage[T1]{fontenc} %% Format of hte resulting PDF file.
\usepackage{fancyhdr} %% Package to create a header on each page.
\usepackage{lastpage} %% Used for "Page X of Y" in the header.
%% For this to work, you have to call pdflatex twice.
\usepackage{enumitem} %% Used to change the style of enumerations (see below).

\usepackage{amssymb} %% Definitions for math symbols.
\usepackage{amsmath} %% Definitions for math symbols.
\usepackage{amsthm}
\usepackage{amsfonts}
\usepackage{mathtools}

\usepackage[english]{babel}

\usepackage{graphicx}
\usepackage{hyperref}
\setlength{\parindent}{0cm}

%\usepackage[left=2.00cm, right=2.00cm, top=2.00cm, bottom=2.00cm ]{geometry}

\newtheorem*{zz*}{z.Z}
\newtheorem*{theorem*}{Theorem}
\newtheorem*{claim*}{claim}
\newtheorem{claim}{claim}
\newtheorem*{prop*}{Propsition}
\newtheorem*{lemma*}{Lemma}
\newtheorem{lemma}{Lemma}

%\title{Lineare Algebra Serie 7}
\date{\today}


\usepackage{tikz}  %% Pagacke to create graphics (graphs, automata, etc.)
\usepackage{tikz-cd} %%commutativ system
\usetikzlibrary{automata} %% Tikz library to draw automata
\usetikzlibrary{arrows}   %% Tikz library for nicer arrow heads


%% Left side of header
\lhead{\course\\\semester\\Serie \homeworkNumber}
%% Right side of header
\rhead{\authorname\\Page \thepage\ of \pageref{LastPage}}
%% Height of header
\usepackage[headheight=36pt]{geometry}
%% Page style that uses the header
\pagestyle{fancy}

\newcommand{\authorname}{Ken Walzer}
\newcommand{\semester}{HS 2023}
\newcommand{\course}{Algebraic number theory}
\newcommand{\homeworkNumber}{2}
\newcommand{\Zn}{\mathcal{Z}}
\newcommand{\Vn}{V_{\mathbb{P}^n}}
\newcommand{\RR}{\mathbb{R}}
\newcommand{\real}{\mathrm{Re}}
\newcommand{\imag}{\mathrm{Im}}
\newcommand{\CC}{\mathbb{C}}
\newcommand{\CCC}{\mathcal{C}}
\newcommand{\QQ}{\mathbb{Q}}
\newcommand{\ZZ}{\mathbb{Z}}
\newcommand{\NN}{\mathbb{N}}
\newcommand{\FF}{\mathbb{F}}
\newcommand{\Mat}{\mathrm{Mat}}
\newcommand{\Aut}{\mathrm{Aut}}
\newcommand{\PP}{\mathbb{P}}
\newcommand{\A}{\mathbb{A}}
\newcommand{\OO}{\mathcal{O}}
\newcommand{\kk}{\mathbf{k}}
\newcommand{\aId}{\mathfrak{a}}
\newcommand{\ii}{\sqrt{-1}}
\newcommand{\GL}{\mathrm{GL}}
\newcommand{\SL}{\mathrm{SL}}
\newcommand{\orb}{\mathrm{orb}}
\newcommand{\stab}{\mathrm{stab}}
\newcommand{\ord}{\mathrm{ord}}
\newcommand{\id}{\mathrm{id}}
\newcommand{\frnr}{\frac{3+4i}{5}}
\linespread{1.2}

\begin{document}
	
	\section*{Aufgabe 1}
	Let $\varphi = \frac{\sqrt{5} + 1}{2}$ embedded into $\QQ$
	\begin{zz*}
		It holds that
		\begin{enumerate}[label = (\roman*)]
			\item $\varphi$ is integer over $\ZZ$
			\item $\varphi^{-1}$ is integer over $\ZZ.$
			\item The ring $\ZZ[\varphi]$ is infinite.
		\end{enumerate}
	\end{zz*}
	\begin{proof}
		\begin{enumerate}[label = (\roman*)]
			\item It holds that
			\[4\cdot \varphi^2 = 5+2\sqrt{5} + 1 = 6+2\sqrt{5}\]
			and
			\[4(4\cdot \varphi^4 - 12\varphi^2 + 9) = 16\cdot \varphi^4 - 48\varphi^2 + 36 = (4\cdot \varphi^2-6)^2 = 4\cdot5\]
			Thus
			\[4\cdot \varphi^4 - 12\varphi^2 + 9 - 5 = 4\cdot \varphi^4-12\varphi^2 +4 = 4(\varphi^4-3\varphi^2+1) = 0\]
			So we get for $P = X^4 - 3X^2 + 1\in \ZZ[X],$ that $P(\varphi) = 0,$ where $P$ is monic. It follows that $\varphi$ is integer over $\ZZ.$
			\item By (i), it follows that
			\[-1 = P(\varphi) - 1= \varphi^4 - 3\varphi^2 = \varphi\cdot \varphi(\varphi^2 - 3)\]
			Thus
			\[1 = (\varphi\cdot \varphi(\varphi^2 - 3))^2 = \varphi\cdot \varphi^3(\varphi^2 - 3)^2\]
			It follows that $\varphi^{-1} = \varphi^3(\varphi^2-3)^2.$ Then we know that the integral elements over a ring form a ring, thus $\varphi^{-1}$ is integral as well.
			\item We define a map $f: \ZZ\to \ZZ[\varphi],$ given by
			\[f: z\mapsto z\cdot \varphi.\]
			We claim that this map is injective. Indeed, let $s,t\in \ZZ,$ such that $f(s) = f(t).$ This gives
			\[s\varphi = t\varphi.\]
			Thus $(s-t)\varphi = 0.$ But $\varphi\in \QQ\setminus\{0\}$ where $\QQ$ a field so in particualr a integral domain. It follows that $s = t$ and this gives us, that $f$ is injective. Since $\ZZ$ is infinte and $Im(f)\subseteq\ZZ[i]$ it follows that $\ZZ[i]$ is infinite as well.
		\end{enumerate}
	\end{proof}
	\pagebreak
	
	\section*{Aufgabe 2}
	We want to determine wheter $\frnr$ is integer over $\ZZ.$
	\bigskip
	
	We denote $\varphi = \frnr.$ Then we define Then it holds that $\varphi^2 =\frac{9 + 24i - 16}{25} = \frac{-7 + 24i}{25}.$ Since $6\varphi = \frac{18+24i}{5},$ we get
	\[5\varphi^2 - 6\varphi = \frac{-7 + 24i}{5} - \frac{18+24i}{5} = \frac{-25}{5} = -5\]
	Thus for $P = 5X^2 - 6X + 5\in \ZZ[X],$ we get that $\varphi$ forms a root of $P.$ In particular, this implies that $\varphi$ is algebraic over $\QQ$ hence we can define a number field $K = \QQ(\varphi).$ Now determining whether $\varphi$ is integral over $\ZZ$ is equal to determining wheter $\varphi\in \ZZ_{K}.$
	
	Since $\varphi, \bar{\varphi}\notin \QQ$ are the only roots of $P,$ and since $\deg(P) = 2,$ we get that $P$ is irreducible over $\QQ.$ It follows that $\frac{1}{5}P$ forms the minimal polynomial of $\varphi$ in $\QQ.$ Since $5,6$ are coprime, it follows that $\frac{1}{5}P\notin \ZZ[X].$ Lemma 1.26 implies that $\varphi\notin \ZZ_{K}.$ thus $\varphi$ can not be integral over $\ZZ.$
	\pagebreak
	
	\section*{Aufgabe 3}
	Let $K/F$ be a finite field extension.
	\subsection*{(i)}
	\begin{zz*}
		$Tr_{K/F}\in End(K),$ where $K$ considered as $F$-vector space.
	\end{zz*}
	\begin{proof}
		Let $x\in K$
	\end{proof}
	\subsection*{(ii)}
	\begin{zz*}
		$N_{K/F}:K\to K$ is mutliplicatve.
	\end{zz*}
	\pagebreak
	
	\section*{Aufgabe 4}
	Let $p\geq 2$ be a prime number. and $\zeta = \exp\left(\frac{2\pi i}{p}\right).$
	\subsection*{(i)}
	\begin{zz*}
		It holds that $\mathcal{B} = (1,\zeta\dots, \zeta^{p-2})$ forms a basis of $\QQ(\zeta)$ as A $\QQ$ vectorspace.
	\end{zz*}
	\begin{proof}
		Let $P$ be the minimal polynomial of $\zeta$ over $\QQ.$ Then we know that $P$ is cyclotomic polynomial $P = \Phi_p.$ Therefore
		\[[\QQ(\zeta):\QQ] = \deg(P) = \varphi(p),\]
		where $\varphi$ the euler's totient function. Since $p$ a prime number this even implies that $[\QQ(\zeta):\QQ] = p-1.$ By an algebra argument, it follows that
		\[\mathcal{B} = (1,\zeta, \dots, \zeta^{p-2})\]
		forms A basis of $\QQ(\zeta)$ over $\QQ.$
	\end{proof}
	\subsection*{(ii)}
	We want to determine $\Mat_{\varphi_{\zeta},\mathcal{B}}$
	It holds that
	\[\Mat_{\varphi_{\zeta},\mathcal{B}} = \left(\begin{array}{cccccc}
		\mid & \mid & \mid & \cdots & \mid & \mid\\
		(\varphi_{\zeta}(1))_{\mathcal{B}} & (\varphi_{\zeta}(\zeta))_{\mathcal{B}} & (\varphi_{\zeta}(\zeta^2))_{\mathcal{B}} & \cdots & (\varphi_{\zeta}(\zeta^{p-3}))_{\mathcal{B}}& (\varphi_{\zeta}(\zeta^{p-2}))_{\mathcal{B}}\\
		\mid & \mid & \mid & \cdots & \mid & \mid
	\end{array}\right)\]
	Then we know that $(\varphi_{\zeta}(1))_{\mathcal{B}} = (\zeta)_{\mathcal{B}}$ and $(\varphi_{\zeta}(\zeta^k))_{\mathcal{B}} = (\zeta^{k+1})_{\mathcal{B}}$ for all $k\in \{1,\dots, p-3\}.$ Then we know that $P = 1+X+\dots+ X^{p-1}\in \QQ[X]$ since $p$ is prime. Since $P(\zeta) = 0,$ it follows that
	\[(\varphi_{\zeta}(\zeta^{p-2}))_{\mathcal{B}} = (\zeta^{p-1})_{\mathcal{B}} = (-1-\zeta-\cdots - \zeta^{p-2})\]
	This implies that
	\[\Mat_{\varphi_{\zeta},\mathcal{B}} = \begin{pmatrix}
		0 & 0 & 0 & \cdots & 0 & -1\\
		1 & 0 & 0 & \cdots & \vdots & -1\\
		0 & 1 & 0 & \cdots & \vdots & -1\\
		0 & 0 & 1 & \ddots & \vdots &-1\\
		0 & 0 & \ddots & \ddots & 0 & -1\\
		0 & \cdots & \cdots & 0 & 1 & -1
	\end{pmatrix}\]
	\subsection*{(iii)}
	We want to determine $Tr_{\QQ(\zeta)/\QQ}(\zeta)$ and $N_{\QQ(\zeta)/\QQ}(\zeta)$
	
	\pagebreak
	
\end{document}
